% !TEX program = lualatex
\documentclass[aspectratio=43,professionalfonts,handout]{beamer}
\usepackage{beamer_template}
\usepackage{enumerate}

\newcommand{\LectureNum}{10}
\newcommand{\LectureTitle}{第3章 フーリエ解析　1節 フーリエ展開}
\newcommand{\TextPages}{p.\,74-90}
\newcommand{\CourseName}{応用数学}
\renewcommand{\TermName}{前期}

\begin{document}

\begin{frame}{第3章 フーリエ解析　1節 フーリエ展開}
  \begin{block}{本節の内容}
  \begin{itemize}
    \item フーリエ級数　（p.\,74--79）
    \item 一般の周期関数のフーリエ展開　（p.\,80--85）
    \item 複素フーリエ級数　（p.\,86--89）
  \end{itemize}
  \end{block}
  \vfill\hfill{\scriptsize \textcolor{gray}{教科書 p.\,74--90}}
\end{frame}

\begin{frame}{フーリエ級数}
  \begin{block}{}
  {\large フーリエ級数}
  \end{block}
  \vfill\hfill{\scriptsize \textcolor{gray}{p.\,74--79}}
\end{frame}

\begin{frame}{定義：フーリエ展開（周期 2π）}
  \begin{block}{}
  $f(t)$ が周期 $2\pi$ の区分的に滑らかな関数であるとき、\\[2pt]
  次の無限級数に展開できる（フーリエ展開という）\\[2pt]
  \[
    f(t) = a_{0}/2 + \sum_{n=1}^\infty (a_{n}\cos(nt) + b_{n}\sin(nt))
  \]
  \end{block}
  \vfill\hfill{\scriptsize \textcolor{gray}{p.\,74--79}}
\end{frame}

\begin{frame}{定理：三角関数の直交性}
  \begin{block}{}
  \begin{align*}
    \int_{-\pi}^{\pi} \cos(mt)\cos(nt) dt &= { 0 (m \neq n), \pi (m = n \neq 0), 2\pi (m = n = 0) } \\
    \int_{-\pi}^{\pi} \sin(mt)\sin(nt) dt &= { 0 (m \neq n), \pi (m = n) } \\
    \int_{-\pi}^{\pi} \cos(mt)\sin(nt) dt &= 0  （任意の m, n）
  \end{align*}
  \end{block}
  \vfill\hfill{\scriptsize \textcolor{gray}{p.\,74--79}}
\end{frame}

\begin{frame}{定理：フーリエ展開の収束定理}
  \begin{block}{}
  \textbf{条件}：
  \begin{itemize}
    \item $f(t)$ は任意の有限区間で区分的に滑らか
    \item 1周期内で区分的に連続
  \end{itemize}
  \textbf{結論}：
  \begin{itemize}
    \item 連続点 $t_{0}$ では : フーリエ級数 $ \to f(t_{0})$
    \item 不連続点 $t_{0}$ では : フーリエ級数 $ \to (1/2)[f(t_{0}+0) + f(t_{0}-0)]$
  \end{itemize}
  \end{block}
  \vfill\hfill{\scriptsize \textcolor{gray}{p.\,74--79}}
\end{frame}

\begin{frame}{例題1　（p.\,78）}
  \begin{exampleblock}{問題}
    $f(x) = \begin{cases} 0 & (-\pi \leq x < 0) \\ x & (0 \leq x < \pi) \end{cases}, f(x+2\pi) = f(x)$ のフーリエ展開を求めよ\\[2pt]
  \end{exampleblock}
\end{frame}

\begin{frame}{例題1　解答}
  \begin{exampleblock}{解答}
  \begin{enumerate}[1.]
    \item[] $c_{0} = (1/2\pi)\int_{-\pi}^{\pi} f(x)dx = (1/2\pi)\int_{0}^\pi x dx = \pi/4$
    \item[] $a_{n} = (1/\pi)\int_{0}^\pi x \cos(nx) dx = -(1/n^{2}\pi)(1-(-1)^n)$
    \item[] $b_{n} = (1/\pi)\int_{0}^\pi x \sin(nx) dx = (-1)^{n+1}/n$
  \end{enumerate}
  \end{exampleblock}
  \begin{alertblock}{答}
    $f(x) = \pi/4 + \sum_{n=1}^\infty [-(1/n^{2}\pi)(1-(-1)^n)\cos(nx) + (-1)^{n+1}/n \cdot \sin(nx)] = \pi/4 + (-2/\pi \cdot \cos x + \sin x) - (1/2)\sin 2x + \ldots$
  \end{alertblock}
\end{frame}

\begin{frame}{例題2　（p.\,80--81）}
  \begin{exampleblock}{問題}
    $f(x) = \begin{cases} 0 & (-1 \leq x < 0) \\ 1 & (0 \leq x < 1) \end{cases}, f(x+2) = f(x)$ のフーリエ展開を求めよ\\[2pt]
  \end{exampleblock}
\end{frame}

\begin{frame}{例題2　解答}
  \begin{exampleblock}{解答}
  \begin{enumerate}[1.]
    \item[] $c_{0} = (1/2)\int_{-1}^{1} f(x)dx = (1/2)\int_{0}^{1} 1 dx = 1/2$
    \item[] $a_{n} = \int_{-1}^{1} f(x)\cos(n\pi x) dx = \int_{0}^1 \cos(n\pi x) dx = 0$
    \item[] $b_{n} = \int_{-1}^{1} f(x)\sin(n\pi x) dx = \int_{0}^1 \sin(n\pi x) dx = (1/n\pi)(1 - (-1)^n)$
  \end{enumerate}
  \end{exampleblock}
  \begin{alertblock}{答}
    $f(x) = 1/2 + \sum_{n=1}^\infty (1/n\pi)(1-(-1)^n) \sin(n\pi x) = 1/2 + (2/\pi)(\sin \pi x + \sin(3\pi x)/3 + \sin(5\pi x)/5 + \ldots)$
  \end{alertblock}
  {\footnotesize \textcolor{gray}{\textit{$N=20, N=50$ での部分和グラフ（ギブス現象）を p.81 に掲載}}}
\end{frame}

\begin{frame}{例題3　（p.\,80）}
  \begin{exampleblock}{問題}
    $f(t) = |t| (-\pi \leq t \leq \pi), f(t+2\pi) = f(t)$ のフーリエ展開を求めよ\\[2pt]
  \end{exampleblock}
\end{frame}

\begin{frame}{例題3　解答}
  \begin{exampleblock}{解答}
  \begin{enumerate}[1.]
    \item[] $f(t)$ は偶関数なので $b_{n} = 0$
    \item[] $a_{0} = (2/\pi) \int_{0}^{\pi} t dt = \pi$
    \item[] $a_{n} = (2/\pi) \int_{0}^{\pi} t \cos(nt) dt$
    \item[] 部分積分 $: = (2/\pi)\begin{cases} [t \sin(nt)/n]_{0}^\pi - (1/n)\int_{0}^\pi \sin(nt) dt \end{cases}$
    \item[] $= (2/\pi) \cdot (1/n^{2})[\cos(n\pi) - 1] = (2/n^{2}\pi)[(-1)^n - 1]$
    \item[] n が奇数のとき $a_{n} = -4/n^{2}\pi, n$ が偶数のとき $a_{n} = 0$
  \end{enumerate}
  \end{exampleblock}
  \begin{alertblock}{答}
    $f(t) = \pi/2 - (4/\pi)(\cos t + \cos(3t)/9 + \cos(5t)/25 + \ldots)$
  \end{alertblock}
\end{frame}

\begin{frame}{一般の周期関数のフーリエ展開}
  \begin{block}{}
  {\large 一般の周期関数のフーリエ展開}
  \end{block}
  \vfill\hfill{\scriptsize \textcolor{gray}{p.\,80--85}}
\end{frame}

\begin{frame}{定義：フーリエ展開（周期 T = 2l）}
  \begin{block}{}
  \[
    f(t) = a_{0}/2 + \sum_{n=1}^\infty (a_{n}\cos(n\pi t/l) + b_{n}\sin(n\pi t/l))
  \]
  \medskip\textbf{フーリエ係数}：
  \begin{align*}
    a_{0} &= (1/l) \int_{-l}^{l} f(t) dt \\
    a_{n} &= (1/l) \int_{-l}^{l} f(t)\cos(n\pi t/l) dt  (n = 1, 2, ...) \\
    b_{n} &= (1/l) \int_{-l}^{l} f(t)\sin(n\pi t/l) dt  (n = 1, 2, ...)
  \end{align*}
  \end{block}
  \vfill\hfill{\scriptsize \textcolor{gray}{p.\,80--85}}
\end{frame}

\begin{frame}{定理：フーリエ余弦展開（偶関数）}
  \begin{block}{}
  \[
    f(t) = a_{0}/2 + \sum_{n=1}^\infty a_{n}\cos(n\pi t/l)
  \]
  \end{block}
  \vfill\hfill{\scriptsize \textcolor{gray}{p.\,80--85}}
\end{frame}

\begin{frame}{定理：フーリエサイン展開（奇関数）}
  \begin{block}{}
  \[
    f(t) = \sum_{n=1}^\infty b_{n}\sin(n\pi t/l)
  \]
  \end{block}
  \vfill\hfill{\scriptsize \textcolor{gray}{p.\,80--85}}
\end{frame}

\begin{frame}{例題3　（p.\,83）}
  \begin{exampleblock}{問題}
    $f(x) = |x| (-1 \leq x < 1), f(x+2) = f(x)$ のフーリエ展開を求めよ\\[2pt]
  \end{exampleblock}
\end{frame}

\begin{frame}{例題3　解答}
  \begin{exampleblock}{解答}
  \begin{enumerate}[1.]
    \item[] $f(x)$ は偶関数なので $b_{n} = 0,$ フーリエ余弦展開を求めればよい
    \item[] $c_{0} = \int_{-1}^{1} |x| dx = 1$
    \item[] $a_{n} = 2\int_{0}^{1} x \cos(n\pi x) dx = -2(1-(-1)^n)/(n^{2}\pi^{2})$
  \end{enumerate}
  \end{exampleblock}
  \begin{alertblock}{答}
    $f(x) = 1/2 - (4/\pi^{2})(\cos \pi x + \cos(3\pi x)/9 + \cos(5\pi x)/25 + \ldots)$
  \end{alertblock}
  {\footnotesize \textcolor{gray}{\textit{$N=2, N=5$ での収束グラフあり（ p.83 ）}}}
\end{frame}

\begin{frame}{例題4　（p.\,84）}
  \begin{exampleblock}{問題}
    $f(x) = x(1-x) (0 \leq x \leq 1), f(-x) = -f(x), f(x+2) = f(x)$ のフーリエ展開を求めよ\\[2pt]
  \end{exampleblock}
\end{frame}

\begin{frame}{例題4　解答}
  \begin{exampleblock}{解答}
  \begin{enumerate}[1.]
    \item[] $f(x)$ は奇関数なので $c_{0} = 0, a_{n} = 0,$ フーリエ正弦展開を求めればよい
    \item[] $b_{n} = 2\int_{0}^{1} x(1-x) \sin(n\pi x) dx = 4(1-(-1)^n)/(n^{3}\pi^{3})$
  \end{enumerate}
  \end{exampleblock}
  \begin{alertblock}{答}
    $f(x) = (8/\pi^{3})(\sin \pi x + \sin(3\pi x)/27 + \sin(5\pi x)/125 + \ldots) = (8/\pi^{3})\sum_{k=0}^\infty \sin((2k+1)\pi x)/(2k+1)^{3}$
  \end{alertblock}
\end{frame}

\begin{frame}{例題5　（p.\,84--86）}
  \begin{exampleblock}{問題}
    次の等式が成り立つことを証明せよ $: 1 - 1/3 + 1/5 - \ldots = \pi/4$\\[2pt]
  \end{exampleblock}
\end{frame}

\begin{frame}{例題5　解答}
  \begin{exampleblock}{解答}
  \begin{enumerate}[1.]
    \item[] 例題 2(1) の関数は $x = 1/2$ で連続だから
    \item[] $f(1/2) = 1/2 + (2/\pi)(\sin(\pi/2) + \sin(3\pi/2)/3 + \sin(5\pi/2)/5 + \ldots)$
    \item[] $= 1/2 + (2/\pi)(1 - 1/3 + 1/5 - \ldots)$
    \item[] $f(1/2) = 1$ なので $: 1 = 1/2 + (2/\pi)(1 - 1/3 + 1/5 - \ldots)$
  \end{enumerate}
  \end{exampleblock}
  \begin{alertblock}{答}
    $1 - 1/3 + 1/5 - \ldots = \pi/4$
  \end{alertblock}
\end{frame}

\begin{frame}{演習問題（p.\,80--85）}
  \begin{exampleblock}{自分で解いてみよう}
  \begin{description}
    \item[問1] 次の等式 (IV) が成り立つことを証明せよ
    $\int_{-\pi}^{\pi} \cos(mx) \cos(nx) dx = { 0  (m \neq n),  \pi  (m = n \neq 0) }$\\
    \item[問2] 次の周期 $2\pi$ の関数 $f(x)$ のフーリエ級数を求めよ
    $f(x) = x  (-\pi \leq x < \pi),  f(x+2\pi) = f(x)$\\
    \item[問3] 次の関数のフーリエ級数を求めよ
    $(1) f(x) = { 1  (-1 \leq x < 0),  0  (0 \leq x < 1) },  f(x+2) = f(x)$\\
    $(2) g(x) = { 2  (-2 \leq x < 0),  4  (0 \leq x < 2) },  g(x+4) = g(x)$\\
    $(3) h(x) = { 2x+1  (-1/2 \leq x < 0),  0  (0 \leq x < 1/2) },  h(x+1) = h(x)$\\
    \item[問4] 次の関数のフーリエ級数を求めよ
    $(1) f(x) = { -1  (-1 \leq x < 0),  1  (0 \leq x < 1) },  f(x+2) = f(x)$\\
    $(2) g(x) = 1 - |x|  (-1 \leq x < 1),  g(x+2) = g(x)$\\
    $(3) h(x) = \cos(\pi x/4)  (-2 \leq x < 2),  h(x+4) = h(x)$\\
    \item[問5] 次の問いに答えよ
    (1) 例題3の関数のフーリエ級数を用いて、次の公式を導け\\
    $1/1^{2} + 1/3^{2} + 1/5^{2} + \ldots = \pi^{2}/8$\\
    (2) 例題4の関数のフーリエ級数を用いて、次の公式を導け\\
    $1/1^{3} - 1/3^{3} + 1/5^{3} - \ldots = \pi^{3}/32$\\
  \end{description}
  \end{exampleblock}
\end{frame}

\begin{frame}{複素フーリエ級数}
  \begin{block}{}
  {\large 複素フーリエ級数}
  \end{block}
  \vfill\hfill{\scriptsize \textcolor{gray}{p.\,86--89}}
\end{frame}

\begin{frame}{定義：複素フーリエ展開}
  \begin{block}{}
  \[
    f(t) = \sum_{n=-\infty}^{\infty} c_{n} e^{int}
  \]
  \end{block}
  \vfill\hfill{\scriptsize \textcolor{gray}{p.\,86--89}}
\end{frame}

\begin{frame}{一般周期関数のフーリエ展開}
  \begin{block}{}
  \[
    f(t) = \sum_{n=-\infty}^{\infty} c_{n} e^{in\pi t/l}
  \]
  \end{block}
\end{frame}

\begin{frame}{例題6　（p.\,88）}
  \begin{exampleblock}{問題}
    次の周期 4 の関数 $f(x)$ の複素フーリエ級数を求めよ\\[2pt]
    $f(x) = x + 1  (-2 \leq x < 2),  f(x+4) = f(x)$\\[2pt]
  \end{exampleblock}
\end{frame}

\begin{frame}{例題6　解答}
  \begin{exampleblock}{解答}
  \begin{enumerate}[1.]
    \item[] $c_{n} = (1/4) \int_{-2}^{2} (x+1) e^{-in\pi x/2} dx$
    \item[] n \neq 0 のとき $: e^{in\pi} = (-1)^n$ を用いて計算
    \item[] $c_{n} = 2(-1)^n / (in\pi)$
    \item[] $n = 0$ のとき $: c_{0} = (1/4)\int_{-2}^{2}(x+1)dx = 1$
  \end{enumerate}
  \end{exampleblock}
  \begin{alertblock}{答}
    $f(x) = 1 + \sum_{n\neq0} 2(-1)^n/(in\pi) \cdot e^{in\pi x/2}$
  \end{alertblock}
\end{frame}

\begin{frame}{演習問題（p.\,86--89）}
  \begin{exampleblock}{自分で解いてみよう}
  \begin{description}
    \item[問6] 次の周期 2 の関数 $f(x), g(x)$ の複素フーリエ級数を求めよ
    $(1) f(x) = { 0  (-1 \leq x < 0),  1  (0 \leq x < 1) },  f(x+2) = f(x)$\\
    $(2) g(x) = { 0  (-1 \leq x < 0),  x  (0 \leq x < 1) },  g(x+2) = g(x)$\\
    \item[問7] 次の周期 4 の関数 $f(x), g(x)$ の複素フーリエ級数を求めよ
    $(1) f(x) = { 2  (-2 \leq x < 0),  0  (0 \leq x < 2) },  f(x+4) = f(x)$\\
    $(2) g(x) = { 0  (-2 \leq x < 0),  2-x  (0 \leq x < 2) },  g(x+4) = g(x)$\\
  \end{description}
  \end{exampleblock}
\end{frame}

\begin{frame}{練習問題1　（p.\,90）}
  \begin{block}{}
  \begin{description}
    \item[問1] 次の周期 2 の関数 $f(x)$ のフーリエ級数を求めよ
    $f(x) = { 1  (-1 \leq x < 0),  1-x  (0 \leq x < 1) },  f(x+2) = f(x)$\\
    \item[問2] グラフが次の図で示されている周期 4 の関数のフーリエ級数を求めよ
    (1) 三角波 $: f(x) = 2 - |x| (-2 \leq x < 2), f(x+4) = f(x)$ （頂点 $y=2,$ 周期 4 ）\\
    (2) のこぎり波 $: f(x) = x (-1 \leq x < 1), f(x+2) = f(x)$ （周期 2 ）\\
    \item[問3] 次の関数について、以下の問いに答えよ
    $f(x) = x^{2}  (-1 \leq x < 1),  f(x+2) = f(x)$\\
    (1) フーリエ級数を求めよ\\
    $(2) x=0$ および $x=1$ とおくことにより、次の公式を導け\\
    $1 - 1/2^{2} + 1/3^{2} - \ldots = \pi^{2}/12$\\
    $1 + 1/2^{2} + 1/3^{2} + \ldots = \pi^{2}/6$\\
    \item[問4] 次の関数 $f(x)$ の複素フーリエ級数を求めよ
    $f(x) = e^x  (-\pi/2 \leq x < \pi/2),  f(x+\pi) = f(x)$\\
  \end{description}
  \end{block}
\end{frame}

\end{document}
